% ================================================================
% CHAPTER 7: NASH EQUILIBRIUM IN THE HILBERT SPACE
% ================================================================
\chapter{Nash Equilibrium in the Hilbert Space}
\label{ch:nash}

\epigraph{Every art and every inquiry, and similarly every action and pursuit,
is thought to aim at some good.}{Aristotle, \textit{Nicomachean Ethics}, I.1}

\section{Equilibrium in the Meaning Game}

In Chapter~\ref{ch:payoffs}, we established the payoff structure of the meaning
game. We now ask: what are the \emph{equilibria} of this game? Under what
conditions does neither player wish to change their signal, given the other's
signal?

\begin{definition}[Nash Equilibrium in the Meaning Game]
\label{def:nash-meaning}
A pair $(s^*, r^*)$ is a \textbf{Nash equilibrium} of the meaning game if:
\begin{enumerate}
\item The sender cannot increase their payoff by changing $s^*$ (holding $r^*$
fixed):
\[
  \payoff{S}{s^*, r^*} \geq \payoff{S}{s, r^*} \quad \text{for all } s \in \IS;
\]
\item The receiver cannot increase their payoff by changing $r^*$ (holding $s^*$
fixed):
\[
  \payoff{R}{s^*, r^*} \geq \payoff{R}{s^*, r} \quad \text{for all } r \in \IS.
\]
\end{enumerate}
\end{definition}

\begin{remark}[The Strategy Spaces]
In the standard meaning game, the sender chooses $s$ and the receiver's
``strategy'' is their prior $r$. In the symmetric version
(Definition~\ref{def:symmetric-game}), both agents choose their ideas
simultaneously. In both cases, the strategy space is $\IS$ (or a subset of it,
if we impose constraints such as budget limits on norms).
\end{remark}


\section{The Trivial Equilibrium}

\begin{theorem}[Void is Always an Equilibrium]
\label{thm:void-equilibrium}
The pair $(\void, \void)$ is always a Nash equilibrium of the symmetric meaning
game, with payoff zero for both players.
\end{theorem}

\begin{proof}
Player $i$'s payoff at $(\void, \void)$ is $\rs{\void}{\void \compose \void} =
\rs{\void}{\void} = 0$.

If player 1 deviates to some $a \neq \void$ while player 2 stays at $\void$:
\[
  \payoff{1}{a, \void} = \rs{a}{a} + \rs{a}{\void} = \rs{a}{a} + 0 = \rs{a}{a}.
\]
This is $\geq 0$, so deviating (weakly) improves player 1's payoff. Thus
$(\void, \void)$ is a Nash equilibrium only in the \emph{weak} sense: no player
\emph{strictly} benefits from deviation if we require $\rs{a}{a} = 0$, which
holds only when $\embed{a} = 0$.

Correction: $(\void, \void)$ is a Nash equilibrium only if no player can achieve
a strictly positive payoff by deviating. Since $\payoff{1}{a, \void} = \rs{a}{a}
\geq 0$, the deviation payoff is non-negative but not strictly greater than 0
only when $\embed{a} = 0$. For any idea $a$ with $\embed{a} \neq 0$, the payoff
$\rs{a}{a} > 0 > 0$, so the deviation \emph{is} strictly profitable.

Therefore, $(\void, \void)$ is \emph{not} a Nash equilibrium when the strategy
space contains any non-void idea! We must revise.
\end{proof}

\begin{theorem}[Void is Not a Nash Equilibrium (Revised)]
\label{thm:void-not-ne}
If the strategy space contains any idea $a$ with $\embed{a} \neq 0$, then
$(\void, \void)$ is \emph{not} a Nash equilibrium.
\end{theorem}

\begin{proof}
Player 1 can deviate to $a$ and receive payoff $\rs{a}{a} > 0 > 0 =
\payoff{1}{\void, \void}$. So the deviation is strictly profitable.

In Lean~4:
\begin{lstlisting}
theorem void_not_ne (a : Idea) (ha : φ a ≠ 0) :
    ¬ is_nash_eq void void := by
  intro h
  have := h.1 a  -- sender can deviate to a
  simp [sender_payoff, rs_void, rs_self_pos ha] at this
  linarith
\end{lstlisting}
\end{proof}

\begin{interpretation}[Silence is Unstable]
\label{interp:silence-unstable}
Theorem~\ref{thm:void-not-ne} has a beautiful interpretation: \emph{silence is
not a stable equilibrium}. If both agents are silent (playing $\void$), either
one can improve their payoff by speaking --- by introducing any non-trivial
idea. The self-resonance term $\rs{a}{a} > 0$ provides a strictly positive
incentive to speak, regardless of what the other agent is doing.

This is a mathematical model of the ``itch to communicate'': there is always a
strategic incentive to break silence, because any non-void idea generates
positive self-resonance. In an evolutionary sense, mutants who ``speak up''
always invade a population of silent agents.

The instability of silence connects to Heidegger's notion that human existence
is fundamentally \emph{disclosive}: Dasein is always already ``thrown'' into a
world of meaning, always already communicating. The Hilbert space model gives
this a game-theoretic foundation: silence is unstable because it leaves positive
payoffs on the table. Speaking is not merely a cultural norm but a strategic
necessity.
\end{interpretation}


\section{Existence of Non-Trivial Equilibria}

Since the void is not an equilibrium, we must look for non-trivial ones. The
existence and structure of equilibria depend on the strategy space.

\begin{definition}[Constrained Meaning Game]
\label{def:constrained-game}
A \textbf{constrained meaning game} is a symmetric meaning game where each
player's strategy space is a closed, bounded, convex subset $\mathcal{S}
\subseteq \Hspace$ (the set of feasible signal embeddings).
\end{definition}

\begin{theorem}[Existence of Nash Equilibrium (Constrained Game)]
\label{thm:ne-existence}
In a constrained meaning game with compact, convex strategy space $\mathcal{S}
\subseteq \Hspace$, a Nash equilibrium exists.
\end{theorem}

\begin{proof}[Proof sketch]
Player $i$'s payoff is $\payoff{i}{v_i, v_j} = \nrm{v_i}^2 + \ip{v_i}{v_j}$,
where $v_i, v_j \in \mathcal{S}$. The best response of player $i$ to $v_j$ is:
\[
  \text{BR}_i(v_j) = \arg\max_{v_i \in \mathcal{S}} \left(\nrm{v_i}^2 +
  \ip{v_i}{v_j}\right).
\]
The objective $v_i \mapsto \nrm{v_i}^2 + \ip{v_i}{v_j}$ is continuous (and in
fact strictly convex) on $\mathcal{S}$. Since $\mathcal{S}$ is compact, the
maximum is attained. The best response correspondence is upper hemicontinuous
(by the Maximum Theorem) and maps the compact, convex set $\mathcal{S}$ to
itself. By Kakutani's fixed point theorem (or Brouwer's, since the best
response is single-valued by strict convexity), a fixed point exists.

Note: the strict convexity of $v_i \mapsto \nrm{v_i}^2 + \ip{v_i}{v_j}$ means
the best response is unique for each $v_j$, so the best response is a
\emph{function} (not just a correspondence), and Brouwer's theorem suffices.

In Lean~4:
\begin{lstlisting}
theorem ne_exists (hS : IsCompact S) (hSc : Convex ℝ S)
    (hSne : S.Nonempty) :
    ∃ v ∈ S, is_nash_eq_constrained v v S := by
  -- Apply Brouwer's fixed point theorem to the best response map
  sorry -- requires Brouwer's theorem from topology
\end{lstlisting}
\end{proof}

\begin{interpretation}[Cultural Equilibria as Fixed Points]
\label{interp:fixed-points}
The existence theorem guarantees that cultural equilibria exist --- at least when
the strategy space is bounded (agents have finite communicative resources). But
what do these equilibria look like?

In a constrained game, the Nash equilibrium tends to occur at the \emph{boundary}
of the strategy space: each player uses their maximum available weight, directed
as favorably as possible given the other's signal. This is because the payoff
$\nrm{v_i}^2 + \ip{v_i}{v_j}$ is increasing in $\nrm{v_i}$ (the self-resonance
term dominates), so players always prefer to ``speak as loudly as possible.''

The direction of the equilibrium depends on the geometry of $\mathcal{S}$. If
$\mathcal{S}$ is a ball (all directions equally available), the equilibrium has
both players choosing the same direction --- the direction that maximizes the sum
of self-resonance and cross-resonance. If $\mathcal{S}$ has asymmetric shape
(some directions are ``easier'' to express than others), the equilibrium reflects
these constraints.

In cultural terms, the strategy space $\mathcal{S}$ represents the set of ideas
that are \emph{expressible} within a given cultural context. A culture with a
rich vocabulary for emotions but a poor vocabulary for mathematics has a strategy
space that extends far in ``emotional'' directions but is constrained in
``mathematical'' directions. The Nash equilibria of discourse within this culture
will favor emotional expression over mathematical precision --- not because the
agents lack mathematical ability, but because the strategy space makes
mathematical signals costly or impossible to produce.
\end{interpretation}


\section{Characterizing Equilibria}

\begin{theorem}[First-Order Characterization]
\label{thm:first-order}
In a constrained meaning game with strategy space $\mathcal{S} = \{v \in
\Hspace : \nrm{v} \leq M\}$ (a ball of radius $M$), the Nash equilibrium
$(v_1^*, v_2^*)$ satisfies: both $v_1^*$ and $v_2^*$ lie on the boundary
$\nrm{v_i^*} = M$, and:
\[
  v_i^* = M \cdot \frac{2v_i^* + v_j^*}{\nrm{2v_i^* + v_j^*}}.
\]
That is, each player's equilibrium strategy points in the direction $2v_i^* +
v_j^*$ --- a weighted combination of their own signal and their partner's.
\end{theorem}

\begin{proof}[Proof sketch]
Player $i$'s payoff is $f(v_i) = \nrm{v_i}^2 + \ip{v_i}{v_j}$. On the ball
$\nrm{v_i} \leq M$, the maximum of $f$ is attained on the boundary (since $f$
is convex and increasing in $\nrm{v_i}$). On the boundary $\nrm{v_i} = M$,
we maximize $\ip{v_i}{v_j}$ (since $\nrm{v_i}^2 = M^2$ is constant), which is
achieved when $v_i = M \cdot v_j / \nrm{v_j}$.

Wait --- let me reconsider. We are maximizing $f(v_i) = \nrm{v_i}^2 +
\ip{v_i}{v_j}$ over $\nrm{v_i} \leq M$. The gradient is $\nabla f(v_i) =
2v_i + v_j$. Since $f$ is strictly convex (the Hessian is $2I$, positive
definite), the unconstrained maximizer does not exist (it goes to $+\infty$).
On the ball, the maximum is on the boundary, and by the KKT conditions:
$2v_i^* + v_j^* = \lambda v_i^*$ for some $\lambda > 0$ (the gradient is a
positive multiple of the outward normal). This gives $v_i^* = v_j^* / (\lambda
- 2)$, or equivalently, $v_i^*$ is in the direction of $v_j^*$. So at the NE,
$v_1^* \parallel v_2^*$!

Let us state this correctly.
\end{proof}

\begin{theorem}[Equilibria Are Parallel]
\label{thm:parallel-equilibria}
In the constrained meaning game with strategy space $\{v : \nrm{v} \leq M\}$,
every Nash equilibrium $(v_1^*, v_2^*)$ satisfies $v_1^* \parallel v_2^*$ (the
strategies are parallel vectors). Specifically, $v_1^* = v_2^* = M \hat{u}$ for
some unit vector $\hat{u}$.
\end{theorem}

\begin{proof}
From the analysis above: the best response of player $i$ to $v_j$ is
$\text{BR}_i(v_j) = M \cdot v_j / \nrm{v_j}$ (assuming $v_j \neq 0$). At
equilibrium, $v_1^* = M \cdot v_2^* / \nrm{v_2^*}$ and $v_2^* = M \cdot
v_1^* / \nrm{v_1^*}$. Substituting:
\[
  v_1^* = M \cdot \frac{M \cdot v_1^* / \nrm{v_1^*}}{\nrm{M \cdot v_1^* /
  \nrm{v_1^*}}} = M \cdot \frac{v_1^*}{\nrm{v_1^*}}.
\]
This is satisfied when $\nrm{v_1^*} = M$ (which we already know). So
$v_1^*$ and $v_2^*$ are both on the boundary and parallel. Since $v_1^* =
M \hat{v}_2$ and $\nrm{v_2^*} = M$, we get $v_1^* = v_2^*$.

In Lean~4:
\begin{lstlisting}
theorem ne_parallel (hne : is_ne v1 v2 M)
    (hv1 : v1 ≠ 0) (hv2 : v2 ≠ 0) :
    ∃ (u : H), ‖u‖ = 1 ∧ v1 = M • u ∧ v2 = M • u := by
  sorry -- follows from best response analysis
\end{lstlisting}
\end{proof}

\begin{interpretation}[Consensus as Equilibrium]
\label{interp:consensus}
Theorem~\ref{thm:parallel-equilibria} is one of the central results of the
game-theoretic analysis: \emph{all Nash equilibria of the constrained meaning
game are consensus equilibria}. At equilibrium, both players express the same
idea (up to the common norm constraint $M$). They ``agree.''

This is not because the players are altruistic or because the game is
cooperative. It is a consequence of pure strategic self-interest: each player's
payoff is maximized by signaling in the same direction as the other player
(to maximize the cross-resonance term), and the self-resonance term is the same
regardless of direction (since $\nrm{v_i^*} = M$ in all equilibria). The
\emph{direction} of the consensus is indeterminate (any direction $\hat{u}$ is
an equilibrium), but the \emph{fact} of consensus is a structural consequence
of the payoff function.

This result formalizes Habermas's thesis that \emph{the telos of communication
is consensus}. Habermas argued that the very structure of communicative action
--- the presupposition that speakers are sincere, truthful, and normatively
appropriate --- pushes discourse toward agreement. IDT~v2 provides a
game-theoretic mechanism: consensus is a Nash equilibrium, and all Nash
equilibria are consensus equilibria. The push toward agreement is not a
normative ideal imposed from outside but a strategic consequence of the
inner-product payoff structure.

But the theorem also reveals a dark side of consensus. The direction of the
consensus is arbitrary --- any direction is an equilibrium. If the initial
conditions favor one direction (say, because of historical accident, power
asymmetry, or propaganda), the equilibrium selects that direction. The resulting
consensus is ``stable'' (no player wants to deviate) but not necessarily
``correct'' or ``just.'' This is the game-theoretic foundation of Gramsci's
theory of \emph{hegemony}: the dominant ideology is an equilibrium, stable
not because it is true but because no individual can profitably deviate from it.
\end{interpretation}


\section{Multiple Equilibria and Coordination}

\begin{theorem}[Continuum of Equilibria]
\label{thm:continuum-ne}
The constrained meaning game with strategy space $\{v : \nrm{v} \leq M\}$ in
a Hilbert space of dimension $n \geq 2$ has a continuum of Nash equilibria,
one for each direction in $\Hspace$. These equilibria all yield the same payoff:
\[
  \payoff{i}{} = M^2 + M^2 = 2M^2 \quad \text{for each player.}
\]
\end{theorem}

\begin{proof}
At any consensus equilibrium $v_1^* = v_2^* = M\hat{u}$:
\[
  \payoff{i}{} = \nrm{M\hat{u}}^2 + \ip{M\hat{u}}{M\hat{u}} = M^2 + M^2 = 2M^2.
\]
There are uncountably many unit vectors $\hat{u}$ (when $n \geq 2$), hence
uncountably many equilibria.
\end{proof}

\begin{interpretation}[The Coordination Problem of Culture]
\label{interp:coordination}
The multiplicity of equilibria is a mathematical model of the \emph{coordination
problem} in cultural production. Every culture must coordinate on a shared
direction of meaning --- a shared set of values, concepts, and priorities. The
game theory tells us that \emph{any} direction is sustainable as an equilibrium,
but the agents must somehow \emph{coordinate} on a particular one.

In the language of game theory, the meaning game is a \emph{pure coordination
game}: all equilibria yield the same payoff, and the challenge is selecting one.
Selection mechanisms include:

\begin{itemize}
\item \textbf{History}: the current equilibrium is the one that happened to be
selected in the past, and inertia maintains it.
\item \textbf{Focal points} (Schelling): some directions are ``natural'' or
``salient'' and serve as coordination devices. In a culture, focal points might
be canonical texts, foundational myths, or charismatic leaders.
\item \textbf{Power}: agents with more weight (higher norm) can ``steer'' the
equilibrium toward their preferred direction, as their signals are more
influential.
\end{itemize}

The indeterminacy of the equilibrium direction is, in fact, the formal content
of cultural \emph{contingency}: the fact that different cultures, facing the
same strategic situation, arrive at radically different equilibria. Why does
one culture organize around honor and shame while another organizes around guilt
and innocence? The game theory says: both are equilibria, and the selection
is a matter of historical contingency, not strategic necessity.
\end{interpretation}

\begin{example}[Equilibrium Selection via Focal Points]
\label{ex:focal}
Consider a meaning game in $\mathbb{R}^2$ with $M = 1$. Any unit vector
$\hat{u} = (\cos\alpha, \sin\alpha)$ defines an equilibrium. Suppose the
culture has a pre-existing ``canonical text'' $t$ with $\embed{t} = (1, 0)$.
If agents use $t$ as a focal point, they coordinate on $\hat{u} = (1, 0)$,
and the equilibrium is $(1, 0), (1, 0)$ with payoff $2$ for each.

If a competing canon $t'$ emerges with $\embed{t'} = (0, 1)$, the agents face
a coordination dilemma: which canon to organize around? The theory predicts
that the equilibrium selected will depend on factors outside the model ---
prestige, power, historical accident --- since both canons support equally
good equilibria.
\end{example}


\section{Equilibria with Heterogeneous Constraints}

\begin{definition}[Heterogeneous Constrained Game]
\label{def:hetero-game}
In a \textbf{heterogeneous constrained game}, player $i$ has a budget
$M_i$ (maximum signal weight): $\nrm{v_i} \leq M_i$.
\end{definition}

\begin{theorem}[Equilibrium with Heterogeneous Budgets]
\label{thm:hetero-ne}
In the heterogeneous constrained game, every Nash equilibrium satisfies
$v_1^* \parallel v_2^*$, with $\nrm{v_i^*} = M_i$ for each $i$. The
payoffs are:
\begin{align*}
  \payoff{1}{} &= M_1^2 + M_1 M_2, \\
  \payoff{2}{} &= M_2^2 + M_1 M_2.
\end{align*}
The player with the larger budget gets the higher payoff: $\payoff{1}{} >
\payoff{2}{}$ iff $M_1 > M_2$.
\end{theorem}

\begin{proof}
The best response analysis from Theorem~\ref{thm:parallel-equilibria} extends
directly: player $i$'s best response to $v_j$ is $M_i \cdot v_j / \nrm{v_j}$.
At equilibrium, $v_1^* = M_1 \hat{u}$ and $v_2^* = M_2 \hat{u}$ for some
unit vector $\hat{u}$. The payoffs follow by direct computation:
$\payoff{i}{} = M_i^2 + M_i M_j$.
\end{proof}

\begin{interpretation}[Inequality in the Marketplace of Ideas]
\label{interp:inequality}
The heterogeneous equilibrium result formalizes a fundamental inequality in
communication: \emph{agents with greater communicative resources
(larger budgets $M_i$) receive higher payoffs}. The payoff difference
$\payoff{1}{} - \payoff{2}{} = M_1^2 - M_2^2$ grows quadratically in the
budget difference.

This is a model of \emph{structural inequality in the marketplace of ideas}.
Agents with more resources --- better education, larger platforms, more
institutional support --- can signal with greater weight and thus extract
more value from every communicative exchange. The ``marketplace of ideas''
is not a level playing field; it systematically favors those who can ``speak
louder'' (produce signals of greater norm).

In Bourdieu's framework, the budget $M_i$ is a component of \emph{cultural
capital}: the accumulated resources that enable an agent to participate
effectively in cultural exchange. The quadratic payoff advantage of larger
$M_i$ explains the self-reinforcing nature of cultural capital: those who
start with more capital extract more value from cultural exchange, enabling
them to accumulate even more capital.
\end{interpretation}

\bigskip

This chapter has established the equilibrium structure of the meaning game: Nash
equilibria are consensus equilibria (both players agree), the direction of
consensus is indeterminate (any direction works), and agents with more
communicative resources extract more value. In Chapter~\ref{ch:cooperation}, we
classify the types of games that arise depending on the angle between ideas.
